
\chapter{Споры и возвратные платежи}\label{ch:disputes}

Когда держатель карты звонит в~банк и~говорит: \enquote{я этого не покупал},
вчерашняя завершённая транзакция перестаёт быть спокойным прошлым. Она
превращается в~спор, где на каждой стадии заново распределяются деньги,
сроки и~бремя доказывания.

\section{Жизнь диспута: пять стадий}\label{sec:disputes-lifecycle}

Получив жалобу, банк не принимает мгновенного решения о~возврате.
Он запускает регламентированный цикл; каждая стадия задаёт свой
срок, набор документов и финансовые последствия~\cite{visa-core-rules,mc-chargeback-guide}.

Сначала нередко возникает предварительный запрос на сведения. В~отрасли
его называют inquiry или retrieval request: эмитент или его процессор
пытается понять, есть ли основания переводить ситуацию в~формальный
спор. Этап не обязателен, но полезен: часть недоразумений снимается ещё
до возвратного платежа.
\definitionnote{Предварительный запрос (inquiry)}{Запрос на сведения до
  формального спора. Его цель: понять, можно ли снять жалобу
без возвратного платежа и~без перераспределения денег по сети.}

Если спор не снят, начинается формальный возвратный платёж без согласия продавца (chargeback):
эмитент обращается к~схеме, схема дебетует эквайера, а~эквайер уже
переносит риск на ТСП. К~спору прикладывается код основания
спора (reason code): именно он показывает, что нужно доказывать и на
чьей стороне лежит основная обязанность по доказательствам~\cite{visa-core-rules}.

Далее следует ответ ТСП на возвратный платёж, который в~отраслевой речи
зовётся повторным представлением (representment). В Visa
официальная терминология ближе к~словарю \enquote{диспут / ответ на
диспут / предарбитраж / арбитраж}, а~в Mastercard
по-прежнему заметна логика циклов возвратного платежа и~подачи дела. Смысл
одинаков: ТСП должно приносить доказательства, а~не
просто возражения.

Если ответ не убедил эмитента, наступает предарбитраж (pre-arbitration):
эмитент либо принимает ответ, либо поднимает спор ещё раз, уже на более
узкой и~дорогой стадии. Если и~это не помогает, схема выносит дело в
арбитраж (arbitration), где решение становится обязательным, а
проигравшая сторона платит по спорной операции и~по процессуальным
сборам~\cite{visa-core-rules,mc-chargeback-guide}.

\begin{figure}[tbp]
  \centering
  \includefiguresvg[width=.85\linewidth]{assets/figures/ch24-disputes/dispute-lifecycle}
  \caption{Пять стадий диспута: от inquiry до арбитража.}\label{fig:dispute-lifecycle}
\end{figure}

Спор выигрывает сторона, которая заранее собрала связную доказательную
базу. Логи, чеки, сведения о~доставке и~согласии клиента должны жить
дольше самого заказа.
\booknote{В~карточных спорах главное не \enquote{быть правым},
  а~заранее хранить нужные доказательства ровно под тот код основания
спора, который применит схема.}

\section{Как Visa и Mastercard ведут спор}\label{sec:disputes-vcr-mc}

В 2018 году Visa перестроила процесс споров на \emph{Visa Claims
Resolution (VCR)}: стороны обязаны раньше приносить доказательства,
а~лишние обмены между эмитентом и эквайером убраны~\cite{visa-vcr-merchants,visa-core-rules}.

Клиент не узнаёт
списание, звонит в~банк, и~банк вместо формального возвратного платежа сначала
показывает ему дополнительные сведения о~покупке. Если продавец даёт
понятное описание товара, адрес доставки или имя получателя, часть
споров исчезает до формального этапа. Именно на этом и~держится
Order Insight: спор пытаются снять ещё до того, как он успеет перейти
в~дорогую процессуальную стадию~\cite{visa-order-insight}.

У Mastercard логика схожа, но язык другой: официальное
руководство говорит о~циклах возвратного платежа, втором предъявлении
(second presentment) и~подаче
дела в~арбитраж~\cite{mc-chargeback-guide}. Архитектура близка, словарь свёрстан по-разному.
Прояснение покупки до формального спора у Mastercard
закрывают продукты Ethoca, включая Consumer
Clarity~\cite{mc-dispute-management}.

Процессуальные различия заметнее всего в~сроках. Visa
последовательно сокращала окно ответа на возвратный платёж: с 45 дней
до 30, затем до 20, а~с 21~июля 2025~года -- до 9 дней для
US/Canada и 18 дней для остальных регионов~\cite{visa-core-rules}.
Mastercard сохраняет 45-дневное окно для ответа эквайера
в~первом цикле возвратного платежа~\cite{mc-chargeback-guide}. У~ТСП
получаются разные SLA: спор по Visa требует
ответа быстрее, чем аналогичный спор по Mastercard. Точный
срок проверяйте по текущей редакции правил.

Исторически 45-дневные окна отвечали операционной реальности
1970-х--80-х: международная почта с~подтверждением доставки
шла несколько недель, банковская обработка добавляла ещё.
Более длинные окна для поздних стадий (предарбитраж, арбитраж)
давали сторонам время собрать доказательства через несколько
часовых поясов. Сокращение Visa до 9/18 дней отражает переход
на цифровую базу: логи, отметки курьерской службы, данные 3DS
доступны немедленно.

На стадии раннего снятия спора механика тоже различается.
Order Insight передаёт детали покупки эмитенту до
формального возвратного платежа: ТСП предоставляет расширенные
данные о~заказе, а~банк использует их при обращении
клиента~\cite{visa-order-insight}.
Ethoca Consumer Clarity работает через инфраструктуру Mastercard:
ТСП предоставляет расширенные данные о~транзакции, и~эмитент
видит их при рассмотрении жалобы~\cite{mc-dispute-management}.
Отдельный продукт Ethoca Alerts уведомляет ТСП о~споре
на ранней стадии: ТСП может провести возврат до
перехода к~формальному возвратному платежу~\cite{mc-dispute-management}.

\begin{figure}[tbp]
  \centering
  \includefiguresvg[width=.8\linewidth]{assets/figures/ch24-disputes/vcr-vs-mc-flow}
  \caption{Два трека раннего снятия спора: Visa Order Insight и Mastercard Ethoca.}\label{fig:vcr-vs-mc-flow}
\end{figure}

Инструменты раннего снятия спора помогают только там, где спор возник
из-за неузнавания покупки или недостатка сведений. При неполученном
товаре, дефектной услуге или нарушенном условии договора спор всё равно
придётся вести по полной схеме.

В~карточном мире важны сам платёж и~качество его описания. Чем яснее
ТСП обозначено в~данных операции, тем меньше у~банка
поводов переводить неузнавание покупки в~формальный спор.

\section{Коды основания спора: четыре категории Visa}\label{sec:disputes-reason-codes}

Visa делит коды основания спора на четыре большие категории:
мошенничество (10.x), ошибки авторизации (11.x), ошибки обработки
(12.x) и потребительские споры (13.x)~\cite{visa-core-rules}.
Эта классификация заранее показывает, кто должен доказывать свою
позицию и~какие доказательства вообще считаются уместными.

Категория 10.x -- мошенничество. Для интернет-торговли определяющим остаётся
код 10.4, \textit{Other Fraud --- Card-Absent Environment}:
держатель карты утверждает, что не участвовал в~интернет-покупке.
У Visa именно эта категория -- главная арена для
Compelling Evidence 3.0.

Категория 11.x -- нарушение правил авторизации.
Транзакция без нужного разрешения, несохранённый корректный ответ
авторизации, продолжение операции вопреки отказу -- спор уходит сюда.

Категория 12.x описывает ошибки обработки -- спор связан
с~некорректной обработкой, а~не с~намерением клиента: сумма не совпала с
авторизацией, одна и~та же операция попала в~клиринг дважды, транзакция
попала в~схему слишком поздно или с~искажёнными данными.

Категория 13.x описывает потребительские споры: товар не получен, услуга
не оказана, подписка не отменена, возврат обещан, но не проведён.
Здесь важнее всего доказательство исполнения обязательства,
а~не технический лог: доставка, акцепт условий, подтверждение пользования,
дата отмены и~момент начисления следующего биллингового цикла.

У Mastercard собственная нумерация кодов основания спора, но логика
та же: мошенничество, авторизация, ошибка обработки, спор держателя
карты~\cite{mc-chargeback-guide}. Разница в~номерах прячет общую
задачу лишь на поверхности: схема сначала определяет класс спора, а~затем
распределяет бремя доказательства.

Коды создают конкретные экономические стимулы. По коду 10.4
(мошенничество в~дистанционном канале) бремя доказательства на
ТСП -- это стимул инвестировать в 3DS, профилирование устройств
и признаки поведенческого совпадения. По кодам~13.x (потребительские
споры) ключевое доказательство -- исполнение обязательства, а~не
технический лог: стимул улучшать подтверждение доставки
и раскрытие условий до транзакции. ТСП с~сильной доказательной
базой по нужному коду получает другое поведение эмитента: при явном
доказательстве у~эмитента меньше оснований держать спор открытым.

\section{Убедительные доказательства: что собирать}\label{sec:disputes-evidence}

Само выражение \enquote{убедительные доказательства} (compelling evidence)
в~споре означает конкретный набор
фактов под конкретный код основания спора~\cite{visa-core-rules}.
То, что помогает по мошенничеству, по неполученному товару не работает.

В~споре о~дистанционной покупке важно показать привязку транзакции
и к~устройству, и~к~поведению клиента. Сюда относятся IP-адрес,
идентификатор или отпечаток устройства,
учётная запись, адрес доставки, адрес электронной почты, номер телефона,
а~при наличии 3-D Secure результат аутентификации. Один IP-адрес
сам по себе слаб; убедительной становится совокупность признаков.

Для спора о~неполучении товара или услуги важнее подтверждение
исполнения: отметка курьерской службы о~вручении, подпись получателя,
технический лог активации цифрового товара, запись о~первом входе в
сервис, протокол оказания услуги. Для спора об отменённой подписке
нужны дата отмены, условия договора и доказательство того, что клиент
видел эти условия до списания.

Документальная гигиена в~карточном бизнесе не должна начинаться
после возвратного платежа. Логи собирайте заранее: IP-адрес, строку
\texttt{User-Agent}, сведения об устройстве, результат AVS и CVV,
результат 3-D Secure, подтверждение доставки и~связку между первой
клиентской транзакцией (CIT) и последующими списаниями. Без этой
связки даже сильная по сути позиция в~споре быстро слабеет.

\section{Compelling Evidence 3.0: зачем нужен NTID}\label{sec:disputes-ce30}

С 15 апреля 2023 года Visa расширила правила доказательства, введя
Compelling Evidence 3.0 (CE~3.0): инструмент для защиты от
спора о~мошенничестве в~дистанционном канале при наличии двух
предшествующих неоспоренных транзакций с~тем же реквизитом и
совпадающими ключевыми признаками клиента. Это ответ Visa на
ситуацию, когда реальный держатель карты оспаривает
собственную покупку так, как будто она была мошеннической~\cite{visa-ce30-merchant-readiness,visa-core-rules}.

Для защиты по CE 3.0 нужны две
предыдущие транзакции того же ТСП без активного сообщения о
мошенничестве (fraud report) и fraud dispute, попадающие в~окно
120--365 дней до даты спора; кроме
того, между спорной и историческими транзакциями должны совпасть как
минимум два ключевых признака клиента, и~одним из них обязан быть
IP-адрес либо идентификатор устройства~\cite{visa-ce30-merchant-readiness}.

На уровне протокола CE~3.0 работает через Visa Resolve Online
(VROL). Когда ТСП получает fraud dispute с reason code 10.4
(\enquote{Other Fraud --- Card-Absent Environment}), система
проверяет, может ли эквайер предоставить данные CE~3.0
автоматически: две предшествующие транзакции от того же ТСП,
каждая без fraud report, для которых совпадает либо IP-адрес,
либо device ID/fingerprint, плюс минимум один дополнительный
элемент из перечня (IP, device ID, user ID, shipping address).
При подтверждённых совпадениях Visa отклоняет диспут (\textit{dispute})
на стадии \textit{pre-dispute} или \textit{dispute}. ТСП не формирует ответ вручную --
сопоставление происходит на стороне Visa. Условие: заранее передавать
device ID и IP в~данных авторизации и хранить историю
неоспоренных транзакций минимум 365~дней~\cite{visa-ce30-merchant-readiness}.

Visa не требует NTID как отдельный критерий CE 3.0. Но в~подписочном
сценарии сохранённый сетевой идентификатор транзакции (Network
Transaction ID, NTID) и аккуратная связка первой клиентской транзакции
(CIT) с~последующими списаниями помогают ТСП восстановить тот самый
исторический след клиента. Вывод следует из сочетания CE 3.0
и Stored Credential Framework~\cite{visa-ce30-merchant-readiness,visa-scf}.

Mastercard развивает First-Party Trust Program: ТСП может передавать
расширенные данные либо в~момент авторизации, либо на стадии спора, а
для защиты ТСП от возвратного платежа программа требует комбинацию
признаков из корзин device factor, delivery factor и additional
identity factors. По терминологии и процессу это самостоятельный инструмент
сопоставимого назначения~\cite{mc-first-party-trust,mc-dispute-management}.

У CE 3.0 узкая область действия: только споры по мошенничеству в
дистанционном канале. Против неполученного товара, дефектной услуги
или ошибочной отмены подписки она не работает.
CE 3.0 помогает там, где спорят о~самом факте участия держателя карты,
но не там, где спорят о~качестве исполнения.

\section{Мониторинговые программы схем}\label{sec:disputes-monitoring}

Системно высокий уровень возвратных платежей грозит попаданием
в~мониторинговые программы схем. Названия и~пороги у~программ
меняются, инвариант один: при слишком высокой доле споров или
мошенничества схема переходит от предупреждений к~плану исправления,
штрафам и, в~крайнем случае, к~прекращению обслуживания~\cite{visa-core-rules,mc-chargeback-guide}.

В~актуальных публичных материалах Visa уже свела мониторинг
споров и мошенничества на уровне эквайера в Visa Acquirer
Monitoring Program (VAMP), а Mastercard по-прежнему описывает
набор отдельных программ контроля, включая Business Risk
Assessment and Mitigation (BRAM), Excessive Chargeback
Program (ECP), Excessive Fraud Merchant (EFM) и Questionable
Merchant Audit Program (QMAP)~\cite{visa-vamp-intro,mc-compliance-programs}.

VAMP -- acquirer-level программа: показатели считаются
по портфелю эквайера, а~не отдельно по каждому ТСП.
Программа объединила прежние Visa Dispute Monitoring Program
(VDMP) и Visa Fraud Monitoring Program (VFMP) в~единый
мониторинг. В~публичном документе Visa, действовавшем
с 1~июня 2025~года, для портфеля эквайера сохранены уровни
Above Standard и Excessive, тогда как для ТСП уровень
Above Standard был отменён -- остался единый глобальный
порог Excessive 220~базисных пунктов (при не менее 1\,500
событий фрода и~диспутов в~месяц), enforcement по которому
начался с 1~октября 2025~года; Visa отдельно объявила снижение
этого порога до 150 базисных пунктов с 1~апреля 2026~года~\cite{visa-vamp-fact-sheet-2025}.
Точный порог проверяйте по дате редакции
программы, а~не пересказывайте как одну постоянную цифру. При превышении
порога участник попадает в~мониторинговую программу и обязан внедрять
меры по снижению риска.

У Mastercard программа мониторинга возвратных платежей ТСП лежит
не в BRAM, а~в ECP. Публичные Security Rules and Procedures
выносят BRAM и ECP в~разные программы, а~для ТСП с
превышением порогов по возвратным платежам используют категории
Excessive Chargeback Merchant (ECM) и High Excessive
Chargeback Merchant (HECM)~\cite{mc-compliance-programs,mc-spme-2025}.
Конкретные числовые пороги Mastercard отсылает к~руководству
Data Integrity Monitoring Program. В~публичном тексте
корректнее говорить об отдельной программе ECP и связанной
эскалации у~эквайера, а~не о~фиксированном пороге BRAM.
Широко используемые ориентиры: для ECM количество возвратных платежей
$\geq$~100 в~месяц при доле возвратных платежей (chargeback ratio) $\geq$~1,5\,\%;
для HECM -- $\geq$~300 при доле $\geq$~3,0\,\%.
Доля возвратных платежей рассчитывается как количество возвратных платежей за текущий
месяц, делённое на количество \textit{e-commerce}-транзакций за
предшествующий месяц~\cite{mc-compliance-programs}.

Схема ожидает от плана ремедиации конкретных мер:
улучшение дескриптора, подключение к Order Insight или Ethoca,
ужесточение проверки 3DS, пересмотр политики возвратов, ограничение
проблемных категорий товаров. ТСП должно показать снижение
показателей в~пределах 90--120 дней, иначе эскалация продолжается.

\begin{figure}[tbp]
  \centering
  \includefiguresvg[width=.65\linewidth]{assets/figures/ch24-disputes/monitoring-escalation}
  \caption{Лестница эскалации мониторинговой программы.}\label{fig:monitoring-escalation}
\end{figure}

Схема смотрит на сочетание
абсолютного объёма и относительной доли. Сезонность, скользящее окно и
перекрёстная миграция споров между календарными месяцами легко искажают
картину. ТСП полезнее считать свои показатели отдельно по
спорам, мошенничеству и возвратам, чем смотреть на один общий дашборд.

Мониторинговая программа подталкивает ТСП своевременно
выстраивать защиту и~избегать массового спора.
Чем лучше дескриптор, чище доставка и~раньше собраны
доказательства, тем реже схема включает жёсткий режим.

\section{Диспуты в~системе \enquote{Мир}}\label{sec:disputes-mir}

Карты \enquote{Мир} обрабатываются ОПКЦ НСПК (глава~\ref{ch:mir}),
и регламент диспутов по ним задают Правила платёжной системы
\enquote{Мир}, а~не Visa Resolve Online или Mastercom~\cite{nspk-mir-rules}.\footnote{Конкретные
коды оснований, сроки и~тарифы стадий диспута публикуются НСПК
в~правилах платёжной системы \enquote{Мир} для участников;
числовые значения и~состав кодов сверяются с~действующей
редакцией перед публикацией.} Архитектурно регламент НСПК
повторяет общую модель карточного диспута: возвратный платёж
(chargeback), повторное представление (representment), претензия
(pre-arbitration), арбитраж. Различия лежат в~кодах оснований,
сроках и~тарифах за стадии.

Структура регламента \enquote{Мир} делит основания диспутов
на несколько групп: операции, оспариваемые держателем карты
(не получил товар или услугу, неавторизованная операция, дубль
операции, неверная сумма); операции, оспариваемые эмитентом
по техническим основаниям (неверная информация о~ТСП, нарушение
правил процессинга, технические сбои); фрод-основания. Сроки
оспаривания у~НСПК свои -- максимальный срок подачи диспута
для большинства оснований составляет 180~дней с~даты операции. Тарифы за стадии (плата
эквайера за принятие возвратного платежа, плата при повторном представлении,
плата за арбитраж) фиксированы в~тарифах ПС~\cite{nspk-mir-tariffs}
и~платятся в~рублях.

Для продавца доказательная база и~сроки
ответа на возвратный платёж по карте \enquote{Мир} из плейбука Visa/Mastercard
заимствовать напрямую нельзя: состав обязательных доказательств и
сроки ответа сверяйте с~правилами \enquote{Мир} через эквайера и~его
процессор. Для эквайера существенно, что мониторинговые программы НСПК
(аналоги VAMP/ECP) -- собственные: пороги доли возвратных платежей и~доли фрода,
после которых ТСП попадает под повышенный контроль или
санкции, устанавливает НСПК.

\section{Диспуты в~СБП: почему их почти нет}\label{sec:disputes-sbp}

Держатель карты, недовольный покупкой, нажимает кнопку спора в
банковском приложении, и карточный возвратный платёж начинает свой цикл.
Плательщик через СБП такой кнопки обычно не видит не потому,
что система \enquote{недоделана}, а~потому, что у~неё нет встроенного
карточного механизма принудительного оспаривания уже исполненной
операции.

В~СБП деньги зачисляются получателю в~режиме реального времени. При
ошибочном переводе банк плательщика может запустить процесс
\enquote{Просьба СБП}, но автоматически вернуть деньги нельзя: банк
обращается к~получателю, потому что для возврата требуется его
согласие, если только иной порядок не вытекает из закона, решения суда
или договора с~банком~\cite{cbr-sbp-faq-mistaken-transfer}.
Для оплаты бизнеса при этом существует обычный возврат через банк
продавца, в~том числе частичный, и~деньги при таком возврате
(refund) поступают покупателю моментально~\cite{sbp-faq-business}.
Встроенного сетевого диспута, похожего на карточный, здесь нет.

У~продавца отсюда преимущество: ему не грозят классические диспутные
штрафы и~длинная цепочка от ответа на возвратный платёж до предарбитража.
У~покупателя -- другое ограничение: если товар не пришёл,
защита строится не на схемном возвратном платеже, а~на возврате, который должен
оформить продавец, на переговорах, претензионной работе или обычном
гражданско-правовом споре. Автоматического списания по инициативе банка
покупателя по правилам схемы здесь нет.

СБП подходит как платёжная инфраструктура для быстрого и окончательного перевода,
но не как замена карточной системе споров. В~продукте, который принимает
и карты, и~СБП, эти два потока обрабатывайте по разным правилам.
Единая машина статусов здесь создаёт лишь ложное ощущение унификации.
